\section{Related work}
\paragraph{Benchmarks and leaderboards.}
MMLU \citep{hendrycks2021measuring} gathers 57 subjects, HELM \citep{liang2022holistic} reports per-subject scores for many models, and the Open LLM Leaderboard ranks models by the mean of six task scores computed with the lm-evaluation-harness \citep{gao2023harness}. These are the leaderboards we use. A different choice of tasks can reorder methods \citep{dehghani2021lottery}; we ask the reverse question, how much one model must change under fixed tasks to top the board.

\paragraph{Aggregation as social choice.}
Several works treat tasks as voters and models as candidates, and propose voting rules to aggregate benchmark results \citep{lanctot2023social,rofin2023votenrank}; the Handbook of Computational Social Choice \citep{brandt2016handbook} covers the rules we use (Borda, plurality, Condorcet-consistent rules). The theory of bribery and of the margin of victory asks how many voters must change their votes to change a winner \citep{faliszewski2009bribery,xia2012margin}. Our problem is the score-valued analogue: the voters are tasks, a change is a raise of one score, and the cost is the points raised. We differ in using real leaderboard matrices, continuous capped score changes and cost measures that suit them. We claim no reduction to, and no complexity result for, bribery or margin of victory: hardness for discrete bribery does not transfer automatically to capped real-valued changes, and for rank rules we rely on a mixed-integer solver that proved optimality or infeasibility on every run reported here.

\paragraph{Gaming and robustness of rankings.}
Test-set leakage and rephrased test data inflate scores \citep{zhou2023dont,yang2023rephrased}; \citet{singh2025illusion} show how private testing and uneven sampling shape the Chatbot Arena leaderboard; \citet{huang2025dropping} show that dropping a very small fraction of preference data can change the top-ranked model in Bradley--Terry rankings. We take the mechanism as given and ask what a score-level intervention costs under each rule.

\paragraph{Strategic developers.}
\citet{chen2026leaderboard} model benchmarking as a Stackelberg game between a benchmark designer and several developers, and prove that current benchmarks induce games with no Nash equilibrium among developers, while tune-before-test gives a unique equilibrium. Ours is a single deviator against fixed rivals on public score matrices, the minimal-cost counterpart and first move of such a game; it shares no theorem with theirs.
